\section{Introduction}\label{sec:intro}

Let $(X,Y)$ be a random pair with $Y>0$ and $X\in\R^p$, and suppose that,
conditionally on $X=x$,
\begin{equation}\label{eq:intro-rv}
 \frac{\Pp(Y>ty\mid X=x)}{\Pp(Y>y\mid X=x)}
 \longrightarrow t^{-1/\gamma(x)},\qquad y\to\infty,\quad t>0,
\end{equation}
for a positive function $\gamma$, the conditional tail index. Estimating
$\gamma$ is the first step toward extreme conditional quantiles and tail
probabilities \citep{deHaanFerreira2006,hill1975}, and it suffers from a
particularly severe curse of dimensionality: local estimation must find
observations near $x$ and then keep only an upper fraction of their
responses. A one-dimensional window of width $h$ combined with a tail
fraction $\alpha$ leaves an effective sample of order $n\alpha h$, not $n$
\citep{daouiaEtAl2013,gardesStupfler2014}. When $p$ is large, even moderate
sample sizes are effectively small.

This paper asks the screening question: among $p$ candidate covariates,
possibly with $p$ growing with $n$, which ones does $\gamma$ actually depend
on? Sure independence screening \citep{fanLv2008} answers such questions by
ranking coordinates through inexpensive marginal statistics and keeping a
moderate number for subsequent joint analysis. Model-free variants exist for
general dependence \citep{liZhongZhu2012} and for fixed conditional quantiles
\citep{heWangHong2013}, but the limiting exponent in \eqref{eq:intro-rv} is a
different target: a covariate can move the $0.95$ quantile of $Y$ without
affecting $\gamma$, and conversely. Our simulations make this distinction
concrete.

The screen we propose rests on a single geometric mechanism. Transform each
continuous covariate to $U_j=F_j(X_j)$ and let $\xi_j(u)$ denote the tail
index of $Y$ given $U_j=u$ alone. Conditioning on one coordinate leaves the
others free to range over the fibre $\{u\}\times[0,1]^{p-1}$, so the
projected distribution is a mixture, and a mixture of heavy tails is governed
by its heaviest component. Under a full-support condition,
\begin{equation}\label{eq:intro-envelope}
 \xi_j(u)=\max_{v\in[0,1]^{p-1}}\gamma(u,v):
\end{equation}
the projected index is the upper envelope of the tail-index surface along the
fibre. If coordinate $j$ is inactive, every fibre reaches the global maximum
$\gmax=\max_x\gamma(x)$ and the envelope is flat at $\gmax$. If coordinate
$j$ is active and some of its fibres miss the global argmax set, the envelope
dips below $\gmax$ there. Averaging the envelope over a trimmed interval
gives a score $\Psi_j$ whose gap $\Delta_j=\gmax-\Psi_j$ is zero for every
inactive coordinate and positive for every detectable active one. Ranking
coordinates by increasing $\Psi_j$ is the screen.

The same geometry states the method's limitation. An active coordinate is
invisible when every one of its fibres still meets the global argmax set;
this happens, for instance, for $\gamma(u_1,u_2)=\gamma_0-(u_1-u_2)^2$,
whose argmax is the diagonal. Detectability is a property of the argmax
projections, and we characterize it exactly rather than assuming a generic
signal-strength condition.

Estimation follows the envelope. For each coordinate we slide a window over
the empirical marginal ranks, apply a Hill statistic to the upper
$\alpha$-fraction of the responses in the window, and average the resulting
profile. Empirical ranks give every coordinate the same deterministic window
geometry and remove the marginal scales. The price is a rank-displacement
term, which we control simultaneously over $p$ coordinates by a
Dvoretzky--Kiefer--Wolfowitz bound; it is negligible when $\log(pn)=o(nh^2)$.
The main result bounds the maximal score error by
\[
 O_{\Pp}\!\left(\sqrt{\frac{\log(pn)}{n\alpha h}}\right)+o(\Delta_{\min}),
\]
where $\Delta_{\min}$ is the weakest population gap and the second term
collects biases that our projected-quantile conditions assume
negligible, and yields sure screening when
$\log(pn)=o(n\alpha h\Delta_{\min}^2)$: high dimension is affordable
exactly when the local extreme count dominates the logarithm of the
dimension. A blockwise implementation costs $O(pn\log n)$ time.

Our construction adapts the projected-tail framework of
\citet{gardesPodgorny2025}, who estimate a low-dimensional central tail-index
subspace by minimizing an average of projected local Hill estimates over
projection matrices, in fixed dimension. We keep their projection semantics
and local Hill estimator, and change what is needed for screening: the
projections are the $p$ coordinates themselves, $p$ may diverge, windows are
built from empirical ranks with deterministic local counts, and uniformity
over a compact matrix class is replaced by union bounds over coordinates and
window states. 

Among screening
proposals, the closest is \citet{yoshidaUmezu2026} screen on covariate-dependent
extreme-value indices through kernel conditional Pickands contrasts in a
large-bandwidth regime. Parametric tail-index regression \citep{wangTsai2009}, extended to
high dimensions with regularization by \citet{sasakiTaoWang2026},
exploits a parametric link when one is available. Tail dimension reduction methods
\citep{gardes2018,aghbalouEtAl2024,bousebataEtAl2023,arbelEtAl2024} target
related but distinct conditional-tail objects. 

Our method is deliberately restricted to the heavy-tailed regime, where the conditional tail index is positive, whereas Yoshida and Umezu~\cite{yoshidaUmezu2026} formulate their screening procedure for a general extreme-value index, allowing positive, zero, and negative values through a conditional Pickands construction. This broader tail-domain coverage is accompanied by a substantially different screening mechanism. Yoshida and Umezu use a large-bandwidth regime, $h\to\infty$, in which the kernel-based conditional Pickands estimator is not intended to estimate the marginal conditional extreme-value-index curve itself, but rather to generate a stable contrast with the unconditional extreme-value index. Their sure-screening result then requires a minimum-signal condition on the resulting population marginal utilities. In contrast, our heavy-tail restriction allows us to identify explicitly the population quantity produced by marginalization: the tail index conditional on a single covariate is the essential supremum of the original conditional tail-index surface over the corresponding fibre, and, under full fibre support, its upper envelope. This projection identity yields an exact characterization of marginal detectability in terms of the projections of the global argmax set: an active coordinate is detectable if and only if a positive-measure set of its fibres fails to intersect that set. Thus, rather than imposing marginal detectability through a generic signal-strength assumption, our approach derives when such a signal exists directly from the geometry of the underlying conditional tail-index function. The two procedures also operate in opposite smoothing regimes: Yoshida and Umezu deliberately increase the bandwidth to construct a stable screening contrast, whereas our empirical-rank local Hill profiles retain localization in the covariate direction in order to estimate the projected tail-index envelope.


Section~\ref{sec:population} develops the population theory.
Section~\ref{sec:estimation} defines the estimator and states the screening
results. Section~\ref{sec:simulation} presents the simulation models, the tuning
study together with a rank-aggregated version of the screen built on a
block of nine neighbouring tunings, and a comparison with competing
screens at $p\in\{500,1000,2000\}$ that covers, in the same tables, the
sensitivity to the dimension and to the quantile level of fixed-quantile
screening. Section~\ref{sec:realdata} applies the aggregated screen to
violent-crime rates in U.S. communities.
Section~\ref{sec:discussion} discusses limitations. Proofs are in
Appendix~\ref{app:proofs}, supported by the auxiliary results of
Appendix~\ref{app:aux}.
